% !TeX root = ../main.tex
\begin{rSection}{Work Experience}

    \begin{rWorkExperienceSubsection}
        {System Engineer (Full-time)}
        {April 2025 --- Present}
        {BREXA Technology Inc.}
        {Tokyo, Japan}
        {
            \item Maintain an AWS-backed bulletin-board web application (Python Flask, MySQL, OpenSearch, Lambda, SQS) used daily by front-office staff for search and content discovery.
            \item Improved search reliability by implementing an HTTP 429 retry policy (one retry, 500\,ms wait), verified with pytest mocks and JMeter load tests.
            \item Monitored post-release production health via CloudWatch Logs and smoke-tested bulletin search flows; traced a recurring 500 error to a timestamp parsing defect and reported root cause to the team for remediation.
            \item Designed and implemented an OpenSearch sort field with \texttt{max(publish\_start, last\_updated)} computed in Lambda (delta-update and reindex batch), updating the Flask app to fix bulletin list ordering without a MySQL migration; presented to a non-technical client, released to staging, and awaiting client acceptance testing.
            \item Built React/TypeScript UI components from Figma designs for a client website project, defining shared component boundaries to reduce duplicate implementation across the team.
            \item Built a Voice RAG proof of concept (React, FastAPI, Twilio, OpenAI Realtime API) for call-centre automation, delivering low-latency speech-to-speech dialogue versus a Bedrock + Polly baseline.
        }
    \end{rWorkExperienceSubsection}

    \begin{rWorkExperienceSubsection}
        {Data Scientist (Internship)}
        {September 2022 --- March 2024}
        {WALC Co.}
        {Tokyo, Japan}
        {
            \item Led development of a Python healthcare chatbot with retrieval-augmented generation (RAG), gathering requirements from medical experts and delivering a secure user-facing application.
            \item Engineered an end-to-end AWS data pipeline and automated document processing workflows, improving response quality by 30\% through systematic data curation.
            \item Built a PyTorch CNN model to classify diabetic retinopathy severity, achieving a quadratic weighted kappa of 0.88 through dataset augmentation and pre-trained architecture experimentation.
        }
    \end{rWorkExperienceSubsection}

    \begin{rWorkExperienceSubsection}
        {Machine Learning Engineer (Internship)}
        {September 2022 --- September 2022}
        {Lightblue Technology Co.}
        {Tokyo, Japan}
        {
            \item Built a PyTorch CNN for multi-class image classification, achieving 91.1\% accuracy through architecture experimentation.
        }
    \end{rWorkExperienceSubsection}

\end{rSection}
